% !TEX root = ../algebra_02.tex

\section{Факторгруппы и прямые произведения}

\subsection{Факторгруппа факторгруппы}

\begin{Theorem}{О факторгруппе факторгруппы}{}
	Пусть $H, P \triangleleft G$ и $H \subset P$. Тогда:
	\begin{enumerate}
		\item $H \triangleleft P$;
		\item $P/H \triangleleft G/H$;
		\item $(G/H)/(P/H) \cong G/P$.
	\end{enumerate}
\end{Theorem}

\begin{proof}
	1. Так как $H \triangleleft G$, то для любых $g \in G$ и $h \in H$ выполнено $ghg^{-1} \in H$. Поскольку $P \subset G$, это условие автоматически выполняется и для всех $g \in P$, что означает $H \triangleleft P$.
	
	2--3. Рассмотрим отображение $\pi_P \colon G \to G/P$, заданное правилом $g \mapsto gP$.
	Так как $H \subset P = \ker \pi_P$, существует индуцированный гомоморфизм:
	\[
	\widetilde{\pi}_P \colon G/H \to G/P, \quad gH \mapsto gP.
	\]
	Очевидно, что $\Im \widetilde{\pi}_P = G/P$. Найдем ядро этого гомоморфизма:
	\[
	\ker \widetilde{\pi}_P = \{ gH \mid gP = eP \} = \{ gH \mid g \in P \} = P/H.
	\]
	Следовательно, ядро является нормальной подгруппой: $P/H \triangleleft G/H$.
	Применяя теорему о гомоморфизме к $\widetilde{\pi}_P$, получаем:
	\[
	(G/H)/(P/H) \cong G/P.
	\]
\end{proof}

\begin{Remark}{}{}
	Упражнение: доказать эту теорему, напрямую построив изоморфизм $G/P \xrightarrow{\sim} (G/H)/(P/H)$.
\end{Remark}

\subsection{Произведение подгрупп}

\begin{Remark}{}{}
	Если $H_1, H_2 \subset G$ --- подгруппы, то их произведение $H_1H_2$ не обязательно является подгруппой.
\end{Remark}

\begin{Example}{Произведение подгрупп не всегда подгруппа}{}
	Рассмотрим группу $G = S_3$ и её подгруппы $H_1 = \langle(1\,2)\rangle$, $H_2 = \langle(1\,3)\rangle$.
	Тогда их произведение состоит из элементов:
	\[
	H_1H_2 = \{ e, (1\,2), (1\,3), (1\,2)(1\,3) \}.
	\]
	Порядок $|H_1H_2| = 4$. По теореме Лагранжа, так как 4 не делит $|S_3| = 6$, множество $H_1H_2$ не является подгруппой.
\end{Example}

\begin{Theorem}{О произведении подгрупп}{}
	Пусть $H \triangleleft G$ и $K < G$. Тогда:
	\begin{enumerate}
		\item $HK = KH < G$;
		\item $H \triangleleft HK$;
		\item $H \cap K \triangleleft K$;
		\item $HK/H \cong K/(H \cap K)$ (запомнить чисто 4-й пункт, остальные нужны для его доказательства).
	\end{enumerate}
\end{Theorem}

\begin{proof}
	1. Очевидно, что $e \in HK$. Так как $H \triangleleft G$, левые и правые смежные классы совпадают:
	\[
	HK = \bigcup_{k \in K} Hk = \bigcup_{k \in K} kH = KH.
	\]
	Проверим замкнутость относительно умножения и взятия обратного:
	\begin{align*}
		(HK)(HK) &= H(KH)K = HHKK = HK, \\
		(HK)^{-1} &= K^{-1}H^{-1} = KH = HK.
	\end{align*}
	Следовательно, $HK$ --- подгруппа в $G$.
	
	2. Условие $H \triangleleft G$ влечет $H \triangleleft HK$, так как $HK \subset G$.
	
	3--4. Рассмотрим отображение $\alpha \colon K \to HK/H$, заданное правилом $k \mapsto kH$. Это гомоморфизм.
	Найдем его образ:
	\[
	\Im \alpha = HK/H.
	\]
	Отображение сюръективно, так как любой элемент смежного класса можно представить как $(kh)H = k(hH) = kH = \alpha(k)$, где $k \in K, h \in H$.
	
	Найдем ядро гомоморфизма $\alpha$:
	\[
	\ker \alpha = \{ k \in K \mid kH = eH \} = \{ k \in K \mid k \in H \} = K \cap H.
	\]
	Отсюда $K \cap H \triangleleft K$. По теореме о гомоморфизме получаем требуемый изоморфизм:
	\[
	K/(K \cap H) \cong HK/H.
	\]
\end{proof}

\begin{Remark}{Аналогия с линейной алгеброй}{}
	Из линейной алгебры известно (в качестве упражнения), что $\dim(V/W) = \dim V - \dim W$.
	Для подпространств выполняется изоморфизм:
	\[
	(W + U)/W \cong U/(W \cap U),
	\]
	что в размерностях дает аналогичную формулу: $\dim(W + U) - \dim W = \dim U - \dim(W \cap U)$.
\end{Remark}

\subsection{Прямое произведение}

\begin{Definition}{Внутреннее прямое произведение}{}
	Пусть $H_1, \dots, H_n < G$. Группа $G$ раскладывается во внутреннее прямое произведение подгрупп $H_1, \dots, H_n$, если:
	\begin{enumerate}
		\item $\forall g \in G \quad \exists! \; h_1 \in H_1, \dots, h_n \in H_n \colon g = h_1 \dots h_n$;
		\item $\forall i \neq j \quad \forall h \in H_i, h' \in H_j \colon hh' = h'h$ (элементы из разных подгрупп коммутируют).
	\end{enumerate}
\end{Definition}

\begin{Proposition}{Критерий внутреннего прямого произведения}{}
	Пусть $H_1, \dots, H_n < G$. Тогда $G$ является внутренним прямым произведением $H_1, \dots, H_n$ тогда и только тогда, когда выполняются условие 2 из определения, а также условия $1'$ и $1''$:
	\begin{itemize}
		\item[$1'$)] $G = H_1 \dots H_n$;
		\item[$1''$)] $\forall i \quad H_i \cap (H_1 \dots \widehat{H}_i \dots H_n) = \{e\}$ (где крышечка означает пропуск $i$-го множителя).
	\end{itemize}
\end{Proposition}

\begin{proof}
	($\Rightarrow$) Условие $1'$ очевидно следует из определения.
	Докажем $1''$. Пусть $h \in H_i \cap (\dots)$. Тогда $h$ можно выразить двояко:
	\[
	h = h_1 \dots \widehat{h}_i \dots h_n \quad \text{и} \quad h = e \dots \underbrace{h}_{i\text{-я позиция}} \dots e.
	\]
	Из единственности разложения (условие 1) следует, что $h_j = e$ при $j \neq i$, и $h = e$.
	
	($\Leftarrow$) Пусть $g = h_1 \dots h_n = h'_1 \dots h'_n$, где $h_i, h'_i \in H_i$. Докажем, что $h_n = h'_n$.
	\begin{align*}
		h_1 \dots h_{n-1} &= h'_1 \dots h'_{n-1} h'_n h_n^{-1} \\
		(h'_{n-1})^{-1} \dots (h'_1)^{-1} h_1 \dots h_{n-1} &= h'_n h_n^{-1}
	\end{align*}
	Элемент слева принадлежит $H_1 \dots H_{n-1}$, а элемент справа принадлежит $H_n$. По условию $1''$, их пересечение тривиально, откуда $h'_n h_n^{-1} = e \implies h_n = h'_n$. Аналогично для остальных множителей.
\end{proof}

\begin{Remark}{Условие нормальности}{}
	Условие $2$ эквивалентно тому, что $H_i \triangleleft G$ для всех $i = 1, \dots, n$.
	
	Действительно, если элементы из разных подгрупп коммутируют, то для $h \in H_i$ и $g = g' h_i g'' \in G$ (где $g' = h_1 \dots h_{i-1}$ и $g'' = h_{i+1} \dots h_n$) имеем:
	\[
	ghg^{-1} = g' h_i g'' h (g'')^{-1} h_i^{-1} (g')^{-1} = \widetilde{h} \in H_i,
	\]
	что означает $H_i \triangleleft G$.
	
	Обратно, если $H_i, H_j \triangleleft G$, рассмотрим коммутатор $[h, h']$:
	\[
	[h, h'] = \underbrace{(h h' h^{-1})}_{\in H_j} (h')^{-1} \in H_j \quad \text{и} \quad [h, h'] = h \underbrace{(h' h^{-1} (h')^{-1})}_{\in H_i} \in H_i.
	\]
	Следовательно, $\widetilde{h} = [h, h'] \in H_i \cap H_j$. Так как $i \neq j$, пересечение тривиально ($\widetilde{h} = e$), откуда $hh' = h'h$.
\end{Remark}

\begin{Example}{Примеры прямых произведений}{}
	\begin{enumerate}
		\item Группа $G = \mathbb{C}^*$ является прямым произведением подгрупп $H_1 = \mathbb{R}^*_+$ и $H_2 = \mathbb{T} = \{ z \in \mathbb{C}^* \mid |z| = 1 \}$.
		\item В группе $G = S_3$ подгруппы $H_1 = A_3$ и $H_2 = \langle(1\,2)\rangle$ порождают $G = H_1 H_2$, но это произведение не является прямым. Это полупрямое произведение: $G = H_1 \rtimes H_2$, так как $H_1 \cap H_2 = \{e\}$ и $H_1 \triangleleft G$, но $H_2$ не нормальна.
	\end{enumerate}
\end{Example}

\subsection{Внешнее прямое произведение}

\begin{Definition}{Внешнее прямое произведение}{}
	Пусть $G_1, \dots, G_n$ --- группы. На их декартовом произведении $G = G_1 \times \dots \times G_n$ введем бинарную операцию:
	\[
	(g_1, \dots, g_n)(g'_1, \dots, g'_n) = (g_1 g'_1, \dots, g_n g'_n).
	\]
\end{Definition}

\begin{Statement}{}{}
	$G$ является группой. Нейтральный элемент: $e_G = (e, \dots, e)$. Обратный элемент: $(g_1, \dots, g_n)^{-1} = (g_1^{-1}, \dots, g_n^{-1})$.
\end{Statement}

Рассмотрим подгруппы:
\[
G'_i = \{ (e, \dots, e, g_i, e, \dots, e) \mid g_i \in G_i \} < G.
\]
Отображение $G_i \to G'_i$, заданное как $g \mapsto (e, \dots, g, \dots, e)$, является изоморфизмом.

\begin{Proposition}{}{}
	Группа $G$ является внутренним прямым произведением своих подгрупп $G'_1, \dots, G'_n$.
\end{Proposition}

\begin{proof}
	Любой $g \in G$ представим в виде:
	\[
	g = (g_1, \dots, g_n) = (g_1, e, \dots) \dots (e, \dots, e, g_n) \in G'_1 \dots G'_n.
	\]
	Коммутирование элементов из $G'_i$ и $G'_j$ (при $i < j$) проверяется напрямую:
	\[
	(e, \dots, g_i, \dots, e)(e, \dots, g'_j, \dots, e) = (e, \dots, g_i, \dots, g'_j, \dots, e) = (e, \dots, g'_j, \dots, e)(e, \dots, g_i, \dots, e).
	\]
\end{proof}

\begin{Proposition}{Критерий изоморфизма прямому произведению}{}
	Пусть $H_1, \dots, H_n < G$. Следующие условия эквивалентны:
	\begin{enumerate}
		\item $G$ --- внутреннее прямое произведение $H_1, \dots, H_n$.
		\item Отображение $\alpha \colon H_1 \times \dots \times H_n \to G$, заданное как $(h_1, \dots, h_n) \mapsto h_1 \dots h_n$, является изоморфизмом.
	\end{enumerate}
\end{Proposition}

\begin{proof}
	Отображение $\alpha$ биективно тогда и только тогда, когда выполнено условие 1 из определения внутреннего произведения.
	
	Если выполнено 2 (изоморфизм), то $\alpha$ --- гомоморфизм. Тогда:
	\[
	\alpha(gg') = \alpha((h_1 h'_1, \dots, h_n h'_n)) = h_1 h'_1 \dots h_n h'_n.
	\]
	При этом $\alpha(g)\alpha(g') = h_1 \dots h_n h'_1 \dots h'_n$. Равенство влечет коммутирование элементов из разных подгрупп.
	
	Обратно, если элементы коммутируют и $\alpha$ биективно, покажем гомоморфность $\alpha$. Для $h \in H_i, h' \in H_j$ ($i \neq j$):
	\begin{align*}
		hh' &= \alpha(e, \dots, h_i, \dots, e)\alpha(e, \dots, h'_j, \dots, e) \\
		&= \alpha((e, \dots, h_i, \dots, e)(e, \dots, h'_j, \dots, e)) \\
		&= \alpha((e, \dots, h_i, \dots, h'_j, \dots, e)) \\
		&= \alpha((e, \dots, h'_j, \dots, e)(e, \dots, h_i, \dots, e)) = h'h.
	\end{align*}
\end{proof}

\begin{Example}{}{}
	\begin{enumerate}
		\item $\mathbb{C}^* \cong \mathbb{R}^*_+ \times \mathbb{T} \cong \mathbb{R} \times (\mathbb{R}/\mathbb{Z})$.
		\item Пусть $m, n \in \mathbb{N}$ и $\gcd(m, n) = 1$. Тогда:
		\[
		\mathbb{Z}/mn\mathbb{Z} \cong (\mathbb{Z}/m\mathbb{Z}) \times (\mathbb{Z}/n\mathbb{Z}).
		\]
		\textit{Доказательство:} Рассмотрим элемент $g = ([1]_m, [1]_n)$. Его порядок равен $mn$, так как $kg = 0 \Leftrightarrow m \mid k$ и $n \mid k \Leftrightarrow mn \mid k$. Значит, он порождает всю группу:
		\[
		\langle ([1]_m, [1]_n) \rangle = (\mathbb{Z}/m\mathbb{Z}) \times (\mathbb{Z}/n\mathbb{Z}) \implies \cong \mathbb{Z}/mn\mathbb{Z}.
		\]
	\end{enumerate}
\end{Example}

\subsection{Конечные абелевы группы}

\begin{Theorem}{Основная теорема о конечных абелевых группах}{}
	Пусть $G$ --- конечная абелева группа. Тогда $G$ изоморфна прямому произведению циклических групп примарных порядков:
	\[
	G \cong \prod \mathbb{Z}/p_i^{n_i}\mathbb{Z},
	\]
	где $p_i$ --- простые числа, $n_i \in \mathbb{N}$. Такое представление единственно с точностью до перестановки сомножителей.
\end{Theorem}

\begin{Example}{Возможные абелевы группы порядка 200}{}
	Пусть $|G| = 200 = 2^3 \cdot 5^2$. Чтобы найти все возможные (с точностью до изоморфизма) абелевы группы такого порядка, переберем разбиения степеней простых множителей.
	
	Варианты для $p=2$ (степень 3):
	\begin{itemize}
		\item $2, 2, 2 \implies (\mathbb{Z}/2\mathbb{Z})^3$
		\item $2^2, 2 \implies (\mathbb{Z}/4\mathbb{Z}) \times (\mathbb{Z}/2\mathbb{Z})$
		\item $2^3 \implies \mathbb{Z}/8\mathbb{Z}$
	\end{itemize}
	
	Варианты для $p=5$ (степень 2):
	\begin{itemize}
		\item $5^2 \implies \mathbb{Z}/25\mathbb{Z}$
		\item $5, 5 \implies (\mathbb{Z}/5\mathbb{Z})^2$
	\end{itemize}
	
	Комбинируя их, получаем список всех возможных групп. Например, при множителе $\mathbb{Z}/25\mathbb{Z}$:
	\begin{align*}
		1) \quad & (\mathbb{Z}/2\mathbb{Z})^3 \times (\mathbb{Z}/25\mathbb{Z}) \\
		2) \quad & (\mathbb{Z}/4\mathbb{Z}) \times (\mathbb{Z}/2\mathbb{Z}) \times (\mathbb{Z}/25\mathbb{Z}) \\
		3) \quad & (\mathbb{Z}/8\mathbb{Z}) \times (\mathbb{Z}/25\mathbb{Z}) \cong \mathbb{Z}/200\mathbb{Z}
	\end{align*}
\end{Example}

\subsection{Конечнопорожденные абелевы группы}

\begin{Definition}{Конечнопорожденная группа}{}
	Группа $G$ называется конечнопорожденной, если существует конечный набор элементов $g_1, \dots, g_s \in G$ такой, что они порождают всю группу: $G = \langle g_1, \dots, g_s \rangle$.
\end{Definition}

\begin{Definition}{Подгруппа кручения (периодическая часть)}{}
	Пусть $G$ --- абелева группа. Множество всех её элементов конечного порядка образует подгруппу $T(G)$, которая называется подгруппой кручения:
	\[
	T(G) = \{ g \in G \mid \operatorname{ord}(g) < \infty \} < G.
	\]
\end{Definition}

\begin{Theorem}{Структура конечнопорожденной абелевой группы}{}
	Пусть $G$ --- конечнопорожденная абелева группа. Тогда в $G$ существует свободная подгруппа $H$ такая, что $G$ является внутренним прямым произведением $H$ и $T(G)$:
	\[
	G = H \times T(G),
	\]
	причем:
	\begin{enumerate}
		\item $T(G)$ --- конечная группа;
		\item $H \cong \mathbb{Z}^r$, где $r \ge 0$.
	\end{enumerate}
	Число $r$ называется рангом группы $G$ и является её инвариантом (не зависит от выбора $H$).
\end{Theorem}

\begin{Example}{Неединственность выбора свободной части $H$}{}
	Рассмотрим группу $G = \mathbb{Z} \times (\mathbb{Z}/2\mathbb{Z})$.
	Её подгруппа кручения определяется однозначно:
	\[
	T(G) = \{ (0, \bar{0}), (0, \bar{1}) \}.
	\]
	В качестве свободной подгруппы $H \cong \mathbb{Z}$ можно взять:
	\[
	H = \{ (a, \bar{0}) \mid a \in \mathbb{Z} \} = \langle (1, \bar{0}) \rangle.
	\]
	Однако этот выбор не уникален. Мы можем взять подгруппу $H' = \langle (1, \bar{1}) \rangle$, которая также изоморфна $\mathbb{Z}$ и удовлетворяет условию $G = H' \times T(G)$. Это показывает, что $H$ в разложении определяется неоднозначно.
\end{Example}

\begin{Example}{Об изоморфизме компонент прямого произведения}{}
	Из изоморфизма прямых произведений $G_1 \times G_2 \cong H_1 \times H_2$ в общем случае не следует изоморфизм отдельных компонент ($G_1 \cong H_1$ и $G_2 \cong H_2$).
	Например:
	\[
	C_3 \times C_2 \cong C_6 \times C_1 \quad (\text{обе изоморфны } C_6),
	\]
	однако $C_3 \not\cong C_6$.
	
	% TODO: на скриншоте приведена таблица подстановок a \mapsto b, b \mapsto -a, -a \mapsto 1/b (?), -b \mapsto -1/a (?). Контекст примера неразборчив, поэтому он опущен для сохранения ясности конспекта.
\end{Example}

\subsection{Мультипликативная группа поля}

\begin{Theorem}{О конечной подгруппе мультипликативной группы поля}{}
	Пусть $K$ --- поле, а $G < K^*$ --- конечная подгруппа его мультипликативной группы (без нуля). Тогда $G$ является циклической группой.
\end{Theorem}

\begin{Remark}{Связь с кольцом вычетов}{}
	Из Основного курса алгебры (ОМА) известно, что мультипликативная группа кольца вычетов $(\mathbb{Z}/n\mathbb{Z})^*$ является циклической тогда и только тогда, когда $n = p^k$, $n = 2p^k$ (где $p$ --- нечетное простое, $k \in \mathbb{N}$), либо $n = 1, 2, 4$.
\end{Remark}

\begin{proof}
	Так как $G$ --- конечная абелева группа, по структурной теореме она изоморфна прямому произведению примарных циклических групп:
	\[
	G \cong \prod_{i=1}^l \mathbb{Z}/p_i^{m_i}\mathbb{Z}.
	\]
	Предположим от противного, что группа $G$ не является циклической. Это означает, что в её разложении среди простых оснований $p_i$ есть одинаковые. Без ограничения общности (WLOG), пусть $p_1 = p_2 = p$.
	
	Перейдем временно к аддитивной записи группы $G'$ (где $G' \cong G$).
	В первой компоненте разложения $\mathbb{Z}/p^{m_1}\mathbb{Z}$ рассмотрим элементы вида $\overline{k \cdot p^{m_1-1}}$, где $0 \le k \le p-1$. Умножение любого такого элемента на $p$ дает $\overline{0}$. Всего таких элементов $p$ штук.
	
	Аналогично, во второй компоненте $\mathbb{Z}/p^{m_2}\mathbb{Z}$ элементы вида $\overline{l \cdot p^{m_2-1}}$, где $0 \le l \le p-1$, при умножении на $p$ дают $\overline{0}$. Их также $p$ штук.
	
	Значит, в прямом произведении элементы вида:
	\[
	(\overline{k \cdot p^{m_1-1}}, \overline{l \cdot p^{m_2-1}}, \overline{0}, \dots, \overline{0})
	\]
	будут давать $\overline{0}$ при умножении на $p$. Общее количество таких элементов составляет не менее $p \times p = p^2$:
	\[
	\left| \{ g' \in G' \mid p \cdot g' = 0 \} \right| \ge p^2.
	\]
	
	Возвращаясь к мультипликативной записи в исходной подгруппе $G < K^*$, мы получаем, что множество элементов $x$, порядок которых делит $p$ (то есть $x^p = 1$), состоит как минимум из $p^2$ элементов:
	\[
	\left| \{ x \in G \mid x^p = 1 \} \right| \ge p^2 \implies \left| \{ x \in K^* \mid x^p = 1 \} \right| \ge p^2.
	\]
	
	Это означает, что многочлен $x^p - 1 = 0$ имеет в поле $K$ не менее $p^2$ корней. Однако мы знаем, что над любым полем многочлен степени $p$ может иметь не более $p$ корней.
	Мы пришли к противоречию ($p^2 \le p$), следовательно, наше предположение неверно. Все простые основания $p_i$ в прямом произведении различны.
	
	Раз все $p_i$ различны, то порядки сомножителей попарно взаимно просты, а значит, прямое произведение изоморфно одной большой циклической группе:
	\[
	G \cong \mathbb{Z} \left/ \left( \prod_{i=1}^l p_i^{m_i} \right) \mathbb{Z} \right..
	\]
	Таким образом, $G$ --- циклическая группа. Теорема доказана.
\end{proof}